\section*{Training 1. A price tag on a standard of living · 30 points}
A consumer has continuous, locally nonsatiated preferences and expenditure
function
\[
 e(p,\bar u)=\bar u\,p_1^{1/4}p_2^{3/4},\qquad p\gg0,\quad\bar u>0.
\]
The relevant optima are unique. For a price increase, express CV, EV and the
loss of consumer surplus as \emph{positive monetary losses}.
\begin{enumerate}[label=(\alph*)]
\item \textbf{6 points.} Use Shephard's lemma to recover both Hicksian
demands. Verify $p\cdot h=e$.
\item \textbf{8 points.} Invert the expenditure function to obtain $v(p,m)$.
Use Roy to recover both Marshallian demands. Check that substituting
$\bar u=v(p,m)$ into h gives x.
\item \textbf{10 points.} Income is $m=16$. Prices change from
$p^0=(1,1)$ to $p^1=(16,1)$. Find the initial and final choices and utilities.
Calculate compensating variation, equivalent variation and the Marshallian
consumer-surplus loss
\[
 L_{CS}=\int_1^{16}x_1((s,1),16)\,ds.
\]
Compare the three monetary amounts.
\item \textbf{6 points.} Instead of restoring initial utility, a scholarship
administrator makes the initial basket just affordable at the new prices.
How much money is added? Compare it with CV and calculate the utility the
consumer can now attain after reoptimising.
\end{enumerate}
